\documentclass{article}
\usepackage{graphicx} % Required for inserting images
\usepackage{amsmath}
\usepackage{algorithmic}
\usepackage{algorithm}
\usepackage[polish]{babel}
\usepackage[utf8]{inputenc}
\usepackage[T1]{fontenc}
\usepackage{float}
\usepackage{listings}
\usepackage[table]{xcolor}
\usepackage{diagbox}
\usepackage{geometry}

\title{Algorytmy Równoległe - Laboratorium 7}
\author{Adam Staniszewski}

\lstset{
    language=python,
    basicstyle=\ttfamily,
    keywordstyle=\color{blue},
    commentstyle=\color{gray},
    stringstyle=\color{orange},
    backgroundcolor=\color{lightgray!20},
    frame=single,
    captionpos=b
}

\begin{document}

\setlength{\parindent}{0pt}

\maketitle

\section{Implementacja sekwencyjna}
Dla porównania z innymi metodami, przygotowałem implementację sekwencyjną.

\vspace{1em}
\textbf{Stałe}
\begin{lstlisting}
G = 6.67430e-11  
EPS = 1e-10  
MASS_SCALAR = 1e30 
\end{lstlisting}

\vspace{1em}
\textbf{Obliczanie przyspieszenia}

\begin{lstlisting}
def compute_forces(stars: list[dict]) -> list[dict]:
    for star in stars:
        star['force'] = np.zeros(3)  
        star['acceleration'] = np.zeros(3)  

    for i, star_i in enumerate(stars):
        for j, star_j in enumerate(stars):
            if i != j: 
                r_vec = star_j['position'] - star_i['position']
                r_mag = np.linalg.norm(r_vec)
                
                force = G * star_i['mass'] * star_j['mass'] 
                * r_vec / (r_mag**3 + EPS)
                star_i['force'] += force

        star_i['acceleration'] = star_i['force'] / star_i['mass']

    return stars
\end{lstlisting}

\vspace{2em}
\textbf{Wyniki}

Dla pięciu gwiazd o danej masie i położeniu:
\begin{lstlisting}
stars = [
    {'position': np.array([1.0, 0.0]), 'mass': MASS_SCALAR},
    {'position': np.array([1.0, 2.0]), 'mass': MASS_SCALAR},
    {'position': np.array([2.0, 4.0]), 'mass': MASS_SCALAR},
    {'position': np.array([3.0, 6.0]), 'mass': MASS_SCALAR},
    {'position': np.array([4.0, 8.0]), 'mass': MASS_SCALAR},
]
\end{lstlisting}

\begin{table}[ht]
    \centering
    \begin{tabular}{|l|r|r|r|}
    \hline
    \rowcolor{orange!40} 
    Gwiazda & Przyspieszenie X & Przyspieszenie Y  \\\hline
    1 & 1.80088671e+18 & 2.29336100e+19 \\\hline
    2 & 8.12539152e+18 & -4.34966964e+17 \\\hline
    3 & 5.40209720e+17 & -8.23998822e+17 \\\hline
    4 & -2.02006859e+18 & -4.56778693e+18 \\\hline
    5 & -8.44641935e+18 & -1.71068573e+19 \\\hline
    \end{tabular}
    \caption{\label{tab:widgets}Przyspieszenia gwiazd dla implementacji sekwencyjnej.}
\end{table}

W celu wyliczenia efektywności, przyspieszenia i metryki Karpa-Flatta przeprowadziłem również eksperyment z użyciem \textbf{4000 gwiazd}, które został wygenerowane przez następujący kod:

\begin{lstlisting}
def generate_stars(num_stars):
    stars = []
    for _ in range(num_stars):
        position = np.random.rand(2) * 10  
        mass = MASS_SCALAR 
        stars.append({'position': position, 'mass': mass})
    return stars
\end{lstlisting}

Algorytm sekwencyjny wykonał wszystkie obliczenia w \textbf{96.46} sekundy.

\section{Zadanie 1}
\begin{lstlisting}
def task_1(stars, rank, size):
    N = len(stars)
    stars_per_process = N // size
    local_stars = stars[rank * stars_per_process: (rank + 1) 
    * stars_per_process]

    transfer_buffer = np.copy(local_stars)
    local_acc = np.zeros_like(local_stars)

    comm = MPI.COMM_WORLD

    for _ in range(size - 1):
        left_rank = (rank - 1) % size
        right_rank = (rank + 1) % size

        comm.send(transfer_buffer, dest=left_rank)
        transfer_buffer = comm.recv(source=right_rank)

        for i, star_i in enumerate(local_stars):
            m_i = star_i['mass']
            r_i = star_i['position']

            for star_j  in transfer_buffer:
                r_j = star_j['position']
                m_j = star_j['mass']

                r_vec = r_j - r_i
                r_mag = np.linalg.norm(r_vec)
                force = G * m_i * m_j * r_vec / (r_mag**3 + EPS)
                local_acc[i] += force / m_i

    return local_acc

\end{lstlisting}

Dla tego samego zestawu gwiazd, co w poprzednim zadaniu, otrzymałem zbliżone wyniki.

\begin{table}[ht]
    \centering
    \begin{tabular}{|l|r|}
    \hline
    \rowcolor{orange!40} 
    Ilość rdzeni & Czas wykonania (s)   \\\hline
    Implementacja sekwencyjna & 96.46  \\\hline
    2 & 32.31  \\\hline
    4 & 31.43  \\\hline
    8 & 35.18  \\\hline
    \end{tabular}
    \caption{\label{tab:widgets}Czas wykonania algorytmu dla 4000 gwiazd przy różnej liczbie rdzeni.}
\end{table}

\begin{figure}[H]
    \centering
    \includegraphics[width=0.6\linewidth]{sp1.png}
\end{figure}

\begin{figure}[H]
    \centering
    \includegraphics[width=0.6\linewidth]{ep1.png}
\end{figure}

\begin{figure}[H]
    \centering
    \includegraphics[width=0.6\linewidth]{kfp1.png}
\end{figure}

\section{Zadanie 2}

\begin{lstlisting}
def task_2(stars, rank, size):
    N = len(stars)
    stars_per_process = N // size
    local_stars = stars[rank * stars_per_process: (rank + 1) 
    * stars_per_process]

    transfer_buffer = np.copy(local_stars)
    local_acc = np.zeros_like(local_stars)

    comm = MPI.COMM_WORLD

    for _ in range(size // 2):
        left_rank = (rank - 1) % size
        right_rank = (rank + 1) % size

        comm.send(transfer_buffer, dest=left_rank)
        transfer_buffer = comm.recv(source=right_rank)
        
        for i, star_i in enumerate(local_stars):
            m_i = star_i['mass']
            r_i = star_i['position']

            for star_j  in transfer_buffer:
                r_j = star_j['position']
                m_j = star_j['mass']

                r_vec = r_j - r_i
                r_mag = np.linalg.norm(r_vec)
                force = G * m_i * m_j * r_vec / (r_mag**3 + EPS)
                local_acc[i] += force / m_i

        transfer_buffer = local_stars.copy()

    return local_acc
\end{lstlisting}

\begin{table}[ht]
    \centering
    \begin{tabular}{|l|r|}
    \hline
    \rowcolor{orange!40} 
    Ilość rdzeni & Czas wykonania (s)   \\\hline
    Implementacja sekwencyjna & 96.46  \\\hline
    2 & 40.85 \\\hline
    4 & 20.69  \\\hline
    8 & 19.27  \\\hline
    \end{tabular}
    \caption{\label{tab:widgets}Czas wykonania algorytmu dla 4000 gwiazd przy różnej liczbie rdzeni.}
\end{table}

\begin{figure}[H]
    \centering
    \includegraphics[width=0.6\linewidth]{sp2.png}
\end{figure}

\begin{figure}[H]
    \centering
    \includegraphics[width=0.6\linewidth]{ep2.png}
\end{figure}

\begin{figure}[H]
    \centering
    \includegraphics[width=1\linewidth]{kpf2.png}
\end{figure}

\section{Zadanie 3}

W tym zadaniu wykonujemy znacznie więcej obliczeń, ponieważ korzystamy zarówno z algorytmu z \textbf{zadania 2}, jak i dokonujemy obliczń przy użyciu metody Eulera.  W związku z tym zmniejszyłem liczbę gwiazd do \textbf{1000}.

\vspace{1em}
Korzystając z metody Eulera, aktualizujemy położenie:

\[
\mathbf{r}(t + \Delta t) = \mathbf{r}(t) + \mathbf{v}(t) \cdot \Delta t
\]

i prędkość:
\[
\mathbf{v}(t + \Delta t) = \mathbf{v}(t) + \mathbf{a}(t) \cdot \Delta t
\]

\begin{lstlisting}
def task_3(stars, rank, size, time_steps, dt):
    N = len(stars)
    stars_per_process = N // size
    local_stars = stars[rank * stars_per_process: (rank + 1) 
    * stars_per_process]

    for star in local_stars:
        star['velocity'] = np.zeros(2)

    comm = MPI.COMM_WORLD

    for _ in range(time_steps):
        local_acc = task_2(stars, rank, size)

        for i, star in enumerate(local_stars):
            # Pozycja
            star['position'] += star['velocity'] * dt
            # Prędkość
            star['velocity'] += local_acc[i] * dt

        all_stars = comm.allgather(local_stars)
        stars = [star for sublist in all_stars for star in sublist]

    return local_stars
\end{lstlisting}

\begin{table}[ht]
    \centering
    \begin{tabular}{|l|r|}
    \hline
    \rowcolor{orange!40} 
    Ilość rdzeni & Czas wykonania (s)   \\\hline
    Implementacja sekwencyjna & 779.56  \\\hline
    2 & 244.64 \\\hline
    4 & 183.43 \\\hline
    8 & 157.95 \\\hline
    \end{tabular}
    \caption{\label{tab:widgets}Czas wykonania algorytmu dla 1000 gwiazd przy różnej liczbie rdzeni.}
\end{table}

\begin{figure}[H]
    \centering
    \includegraphics[width=0.6\linewidth]{wyniki.png}
    \caption{Wyniki dla kilku przykładowych gwiazd}
\end{figure}

\begin{figure}[H]
    \centering
    \includegraphics[width=0.6\linewidth]{sp3.png}
\end{figure}

\begin{figure}[H]
    \centering
    \includegraphics[width=0.6\linewidth]{ep3.png}
    \label{fig:enter-label}
\end{figure}

\begin{figure}[H]
    \centering
    \includegraphics[width=1\linewidth]{kfp3.png}
\end{figure}

\end{document}
